\chapter{Raggio nucleare}

Come per il caso del raggio atomico, anche il raggio nucleare\mkeyword{raggio nucleare} non ha una definzione precisa e natuale e per questo, come per il raggio atomico, sarà necessario introdurre una definizione operativa. 

Come vedremo tra poco, la densità dei nucleoni e il potenziale nucleare hanno una dipendenza spaziale simile: relativamente costanti su distanze piccole oltre le quali cadono drasticamente a zero. \\
\'E quindi intuitivo descrivere la forma nucleare con due parametri: il raggio medio\mkeyword{raggio medio}, in corrispondenza del quale la densità è pari alla metà del suo valore centrale, e lo \emph{skin thickness}\mkeyword{skin thickness}, cioè l'intervallo di spazio in cui la densità scende da un valore prossimo al suo massimo (il $90\%$) a uno prossimo al suo minimo (il $10\%$).

I problemi che incontriamo derivano dalla difficoltà di stabilire con precisione cosa descriva esattamente la distribuzione; il raggio che misuriamo dipende dal tipo di esperimento condotto per determinare la forma del nucleo. In alcuni esperimenti, come nel caso dello scattering di elettroni ad alta energia, misuriamo l'interazione Coulombiana di una particella carica con i nucleoni. Questi esperimenti determinerebbero quindi la \emph{distribuzione della carica nucleare} (principalmente la distribuzione dei protoni, ma che coinvolge in una certa misura anche la distribuzione dei neutroni, a causa della loro struttura interna).

Prima di affrontare il problema conviene ricordare da dove siamo partiti. Nei capitoli precedenti il raggio nucleare è comparso due volte, sempre in modo indiretto: nel termine coulombiano della formula semiempirica di massa, dove entra attraverso $a_C = \frac{3}{5}\frac{\alpha\hbar c}{r_0}$, e nella teoria di Gamow del decadimento $\alpha$, dove fissa il punto di partenza della barriera. Entrambe le vie portano a
\begin{equation}
    R = r_0 A^{1/3} \qquad , \qquad r_0 \simeq 1{,}25 \ \mathrm{fm}
    \label{eq:raggio_A13}
\end{equation}
Si tratta però di stime \emph{dedotte} da modelli, non di misure dirette. Quello che ci proponiamo ora è di misurare il raggio nucleare, e più in generale la distribuzione di carica, bombardando il nucleo e osservando cosa viene diffuso.

\section{Distrubuzione della carica nucleare}

Un metodo utile per determinare la grandezza e la forma di un oggetto è esaminare la radiazione da esso diffusa. 

Per vedere l'oggetto e i suoi dettagli, la lunghezza d'onda della radiazione deve essere più piccola delle dimensioni dell'oggetto, altrimenti l'effetto della diffrazione oscurerebbe parzialmente o completamente l'immagine. \\
Per i nuclei, il cui diametro è di circa $10 \ \text{fm}$, richiediamo che $\lambda \leq  10 \ \text{fm}$ e quindi, dalla relazione\footnote{Dove abbiamo supposto che $p \approx E/c$ in quanto $pc = \sqrt{E^2 - m^2c^4}$ con $mc^2 \approx 0.5 \ \text{MeV} \ll E$.}
\begin{equation}
    \lambda = \frac{2\pi \hbar}{p} \approx \frac{2\pi \hbar c}{E}
\end{equation}
troviamo che gli elettroni usati per la radiazione devono avere un'energia maggiore di $100 \ \text{MeV}$. 

Fasci di elettroni con energia compresa tra $100 \ \text{MeV}$ e $1 \ \text{GeV}$ possono essere prodotti con un acceleratore ad alta energia e possono essere analizzati con precisione da uno spettrometro per selezionare solo gli elettroni che sono stati diffusi elasticamente dal bersaglio nucleare selezionato.

Adesso, però, dobbiamo rendere il probema quantitativo. La funzione d'onda di un elettrone può essere esperra come un'onda piana di impulso ben definito,
\begin{equation}
    \psi_e \propto e^{i(\mathbf{k}\cdot\mathbf{r} - Et/\hbar)} \qquad , \qquad k = \frac{p}{\hbar}
    \label{eq:onda_piana}
\end{equation}
Rigorosamente un elettrone del fascio è un pacchetto d'onda, ma essendo il fascio monocromatico in impulso tale pacchetto è così stretto in $k$, e di conseguenza così esteso nello spazio, da essere indistinguibile da un'onda piana nella regione del bersaglio.\\

A questo punto il problema è formalmente identico a quello di un'onda che incide su un reticolo di diffusori, e lo affronteremo nello stesso modo: calcoleremo l'ampiezza diffusa da un centro puntiforme, poi da due, poi da una distribuzione continua di carica.

\subsection{Scattering su bersagli puntiformi}

Iniziamo con il caso più semplice, quello della diffusione di una radiazione su bersagli puntiformi.

\begin{figure}[H]
    \centering
    \begin{tikzpicture}[>={Stealth[length=2mm]},scale=0.95,
      every node/.style={font=\small}]
    
    % ---------- onda piana incidente ----------
    \foreach \x in {-2.8,-2.4,-2.0}{
      \draw[thin] (\x,-0.55) -- (\x,2.35);
    }
    \draw[->,thick] (-2.70,2.75) -- (-1.70,2.75) node[midway,above] {$\mathbf{k}_0$};
    \node[align=center] at (-2.4,-1.05) {onda piana\\incidente};
    
    % ---------- asse z ----------
    \draw[dashed] (-3.0,0) -- (7.3,0) node[right] {$z$};
    
    % ---------- origine ----------
    \fill (0,0) circle (1.6pt);
    \node[below left] at (0,0) {$O$};
    
    % ---------- diffusori ----------
    \coordinate (P1) at (0.45,1.60);
    \coordinate (P2) at (1.75,1.10);
    \coordinate (D)  at (6.12,2.47);
    
    \draw[->] (0,0) -- (P1) node[pos=0.62,above left,inner sep=1pt] {$\mathbf{r}_1$};
    \draw[->] (0,0) -- (P2) node[pos=0.50,above left,inner sep=1pt] {$\mathbf{r}_2$};
    \draw[->] (P1) -- (P2) node[pos=0.50,above right,inner sep=1pt] {$\mathbf{d}$};
    
    \fill (P1) circle (2.2pt) node[above,inner sep=2.5pt] {$1$};
    \fill (P2) circle (2.2pt) node[below right,inner sep=1.5pt] {$2$};
    
    % ---------- angolo theta ----------
    \draw[thin] (1.15,0) arc (0:22:1.15);
    \node at (11:1.42) {$\theta$};
    
    % ---------- distanze al rivelatore ----------
    \draw (P1) -- (D) node[pos=0.45,above,inner sep=2.5pt] {$\left|\mathbf{R}-\mathbf{r}_1\right|$};
    \draw (P2) -- (D) node[pos=0.15,above,inner sep=2.5pt] {$\left|\mathbf{R}-\mathbf{r}_2\right|$};
    \draw[dashed] (0,0) -- (D) node[pos=0.78,below,inner sep=3pt] {$R$};
    
    % versore n
    \draw[->,thick] (2.90,1.17) -- (3.70,1.50) node[pos=0.5,below,inner sep=2.5pt] {$\hat{\mathbf{n}}$};
    
    % ---------- rivelatore ----------
    \draw[line width=1.6pt] (6.28,2.08) -- (5.96,2.86);
    \node[above] at (6.05,2.95) {rivelatore};
    
    % ---------- nota ----------
    \node at (4.45,0.55) {$R\gg r_1,\,r_2$};
    
    \end{tikzpicture}
    \caption{Scattering di due corpi puntiformi.}
    \label{fig:scena_due}
\end{figure}

Prima di iniziare ad analizzare lo scattering vero e proprio, però, è conveniete trovare una relazione tra la sezione d'urto dfferenziale e l'ampiezza di scattering $A(\theta, \phi)$. \\
Lontano dal centro di scattering, la funzione d'onda si può scrivere come 
\begin{equation}
\psi(\mathbf{r}) \simeq \underbrace{e^{i k z}}_{\text {onda incidente }}+\underbrace{A(\theta, \phi) \frac{e^{i k r}}{r}}_{\text {onda diffusa }}
\end{equation}
Si nota che ci sono due contributi: il primo, $e^{ikz}$, rappresenta un'onda piana incidente che si propaga lungo $z$, mentre il secondo, $A(\theta, \phi) e^{ikr}/r$, un'onda uscente sferica, che descrive le particelle diffuse nelle diverse direzioni. La funzione $A(\theta, \phi)$ determina quanto è intensa l'onda diffusa in una certa direzione.

Per lo scattering dobbiamo sapere quante particelle attraversano una superficie per unità di tempo e questa proprietà è descritta dalla densità di corrente di probabilità 
\begin{equation}
    \mathbf{j} = \frac{\hbar}{m}\operatorname{Im}(\psi^* \nabla \psi )
\end{equation}
Calcoliamo separatamente la corrente incidente e quella diffusa. 

Per l'onda incidente si ha $\psi_i = e^{ikz}$ e $\nabla \psi_i = ik\hat{\mathbf{z}} e^{ikz}$, quindi 
\begin{equation}
\label{eq:corrente_incidente}
    \mathbf{j}_i = \frac{\hbar}{m} \operatorname{Im}(e^{-ikz} i k \hat{\mathbf{z}}e^{ikz}) = \frac{\hbar k}{m} \hat{\mathbf{z}}
\end{equation}

Per l'onda diffusa si ha $\psi_s = A(\theta, \phi) e^{ikr}/r$ e $\partial_r \psi_s \approx ik A(\theta, \phi) e^{ikr}/r$, quindi 
\begin{equation}
\label{eq:corrente_diffratta}
    j_{s,r} = \frac{\hbar}{m} \operatorname{Im}\left( \psi_s^* \frac{\partial \psi_s}{\partial r} \right) = \frac{\hbar k}{m} \frac{|A(\theta,\phi)|^2}{r^2}
\end{equation}

Per definizione, la sezione d'urto differenziale è il rapporto tra il numero di particelle diffuse per unità di tempo nell'angolo solito $d\Omega$ e il flusso incidente per unità di superficie. \\
Il numero di particelle diffuse al secondo attraverso una superficie $dS$ è proporzionale al prodotto tra $j_{s,r}$ e la supeficie stessa, mentre il flusso incidente è proporzionale a $j_i$, da cui si ottiene 
\begin{equation}
    d\sigma = \frac{j_{s,r} dS}{j_i} = \frac{j_{s,r} r^2 d\Omega}{j_i} \implies \frac{d\sigma}{d\Omega} = \frac{j_{s,r} r^2}{j_i}
\end{equation}
Dove abbiamo usato la relazione $dS = r^2 d\Omega$ in quanto supponiamo di considerare una superficie sferica. Sostituendo nell'equazione sopra le equazioni \eqref{eq:corrente_incidente} e \eqref{eq:corrente_diffratta} si ottiene 
\begin{tcolorbox}[colback=black!5!white, colframe=black!45!white, boxrule=0.8pt]
\begin{equation}
    \frac{d\sigma}{d\Omega} = |A(\theta, \phi)|^2
\end{equation}
\end{tcolorbox}

Adesso possiamo iniziare la vera e propria analisi dello scattering. 

L'onda incidente è piana, cioè il fascio è monocromatico in impulso. I due diffusori sono \emph{identici}, e questo garantisce che diffondano con la stessa ampiezza, che indicheremo con $f(\theta,\phi)$ per distinguerla dall'ampiezza $A(\theta,\phi)$ del bersaglio complessivo. Supponiamo inoltre che la diffusione sia \emph{elastica}, per cui $|\mathbf{k}| = |\mathbf{k}_0| = k$, e che avvenga per \emph{singolo urto}, cioè che l'onda diffusa da un centro non venga a sua volta diffusa dall'altro: è quest'ultima ipotesi a rendere il problema lineare e risolubile esattamente.\\

I due centri si trovano nelle posizioni $\mathbf{r}_1$ e $\mathbf{r}_2$ e il rivelatore è posto in $\mathbf{R}$, lontano dalla zona di interazione, nella direzione $\hat{\mathbf{n}} = \mathbf{R}/R$. Chiamiamo $\mathbf{k} = k\,\hat{\mathbf{n}}$ il vettore d'onda uscente.

I due diffusori non emettono in fase, e per due motivi distinti. Il primo è che l'onda incidente non li investe nello stesso istante: arrivando in $\mathbf{r}_j$ essa ha accumulato la fase
%
\begin{equation}
    \psi_{\text{inc}}(\mathbf{r}_j) = e^{i\mathbf{k}_0\cdot\mathbf{r}_j}
    \label{eq:fase_incidente}
\end{equation}
Il secondo è che i due centri non distano ugualmente dal rivelatore. Il contributo del diffusore $j$ nel punto $\mathbf{R}$ è dunque
%
\begin{equation}
    \psi_j(\mathbf{R}) = e^{i\mathbf{k}_0\cdot\mathbf{r}_j}\; f(\theta,\phi)\; \frac{e^{ik\left|\mathbf{R}-\mathbf{r}_j\right|}}{\left|\mathbf{R}-\mathbf{r}_j\right|}
    \label{eq:contributo_j}
\end{equation}

Poiché $R \gg r_j$ sviluppiamo la distanza diffusore-rivelatore,
%
\begin{equation*}
    \left|\mathbf{R}-\mathbf{r}_j\right| = R\sqrt{1 - \frac{2\,\hat{\mathbf{n}}\cdot\mathbf{r}_j}{R} + \frac{r_j^2}{R^2}} \;\simeq\; R - \hat{\mathbf{n}}\cdot\mathbf{r}_j
\end{equation*}
Questo sviluppo va però trattato con due criteri diversi a seconda di dove compare. Nell'\emph{esponente} il termine $\hat{\mathbf{n}}\cdot\mathbf{r}_j$ è moltiplicato per $k$, e il prodotto $k\,r_j$ non è affatto piccolo: la lunghezza d'onda è stata scelta proprio confrontabile con le distanze in gioco, ed è questo il motivo per cui l'esperimento è sensibile alla geometria. Al \emph{denominatore} invece lo stesso termine va confrontato con $R$, e $r_j/R$ è trascurabile. Si ottiene quindi
%
\begin{equation}
    \frac{e^{ik\left|\mathbf{R}-\mathbf{r}_j\right|}}{\left|\mathbf{R}-\mathbf{r}_j\right|} \;\simeq\; \frac{e^{ikR}}{R}\,e^{-i\mathbf{k}\cdot\mathbf{r}_j}
    \label{eq:campo_lontano}
\end{equation}

Sostituendo la \eqref{eq:campo_lontano} nella \eqref{eq:contributo_j} le due fasi si combinano in una sola,
\begin{equation}
    \psi_j(\mathbf{R}) \simeq f(\theta,\phi)\,\frac{e^{ikR}}{R}\;e^{i\left(\mathbf{k}_0-\mathbf{k}\right)\cdot\mathbf{r}_j} = f(\theta,\phi)\,\frac{e^{ikR}}{R}\;e^{i\mathbf{q}\cdot\mathbf{r}_j}
    \label{eq:contributo_finale}
\end{equation}
dove si è introdotto il trasferimento di impulso\mkeyword{trasferimento di impulso}
\begin{equation}
    \mathbf{q} \equiv \mathbf{k}_0 - \mathbf{k}
    \label{eq:q_def}
\end{equation}

Essendo l'equazione d'onda lineare, i due contributi si sommano a livello di ampiezza e non di intensità,
\begin{equation}
    \psi_{\text{diff}}(\mathbf{R}) = \psi_1 + \psi_2 = \underbrace{f(\theta,\phi)\left(e^{i\mathbf{q}\cdot\mathbf{r}_1} + e^{i\mathbf{q}\cdot\mathbf{r}_2}\right)}_{\textstyle A(\theta,\phi)}\;\frac{e^{ikR}}{R}
    \label{eq:sovrapposizione}
\end{equation}
La quantità raccolta è proprio l'ampiezza totale del bersaglio, cioè il coefficiente dell'onda sferica uscente: il sistema dei due diffusori si comporta, visto da lontano, come un unico centro di ampiezza $A(\theta,\phi)$. Applicando allora il risultato già ottenuto, $d\sigma/d\Omega = |A(\theta,\phi)|^2$, si ha
\begin{equation}
    \frac{d\sigma}{d\Omega} = \left|f(\theta,\phi)\right|^2 \left|e^{i\mathbf{q}\cdot\mathbf{r}_1} + e^{i\mathbf{q}\cdot\mathbf{r}_2}\right|^2
    \label{eq:sezione_due}
\end{equation}

Raccogliendo $e^{i\mathbf{q}\cdot\mathbf{r}_1}$, che ha modulo unitario e quindi scompare, e ponendo $\mathbf{d} \equiv \mathbf{r}_2 - \mathbf{r}_1$,
\begin{equation*}
    \left|1 + e^{i\mathbf{q}\cdot\mathbf{d}}\right|^2 = \left(1+e^{i\mathbf{q}\cdot\mathbf{d}}\right)\left(1+e^{-i\mathbf{q}\cdot\mathbf{d}}\right) = 2 + 2\cos\left(\mathbf{q}\cdot\mathbf{d}\right) = 4\cos^2\!\left(\frac{\mathbf{q}\cdot\mathbf{d}}{2}\right)
\end{equation*}
e quindi
\begin{tcolorbox}[colback=black!5!white, colframe=black!45!white, boxrule=0.8pt]
\begin{equation}
    \frac{d\sigma}{d\Omega}\bigg|_{2} = 4\cos^2\!\left(\frac{\mathbf{q}\cdot\mathbf{d}}{2}\right) \ |f(\theta,\phi)|^2
    \label{eq:interferenza_due}
\end{equation}
\end{tcolorbox}

Notando che $|f(\theta,\phi)|^2$ è la sezione d'urto differenziale di un singolo diffusore, possiamo scrivere 
\begin{equation}
    \frac{d\sigma}{d\Omega}\bigg|_{2} = 4\cos^2\!\left(\frac{\mathbf{q}\cdot\mathbf{d}}{2}\right) \ \frac{d\sigma}{d\Omega}\bigg|_{1}
\end{equation}

Il risultato dipende soltanto dalla posizione relativa dei due centri e non dalla scelta dell'origine.

\subsection{Scattering su bersagli continui}

Andiamo adesso a generalizzare quanto trovato nella sezione precedente. 

\'E facile notare che se consideriamo un insieme discreto di cariche la \eqref{eq:sezione_due} diventa 
\begin{equation}
    \frac{d\sigma}{d\Omega} = |f(\theta,\phi)|^2 \left| \sum_{i=1}^Z e^{i \mathbf{q}\cdot \mathbf{r}_i} \right|^2
    \label{eq:somma_discreta}
\end{equation}
dove $f(\theta,\phi)$ è l'ampiezza di diffusione di una singola carica $e$.

La distribuzione di questo insieme discreto, però, può essere rappresentata come una somma di delta di Dirac tridimensionali
\begin{equation}
    n(\mathbf{r}) = \frac{\sum_{i=1}^Z \delta^3(\mathbf{r} - \mathbf{r}_i)}{Z}
\end{equation}
L'introduzione di questa ``densità'' di distribuzione mi permette di passare dalla sommatoria all'integrale, che può essere visto come una generalizzazione ulteriore che mi permette di passare dal caso discreto a quello continuo
\begin{equation}
    \sum_{i=1}^Z e^{i \mathbf{q}\cdot \mathbf{r}_i} = Z \int_{V}n(\mathbf{r}) e^{i\mathbf{q}\cdot\mathbf{r}} d^3r
\end{equation}
dove il fattore $Z$ compare perché $n$ è normalizzata a uno. Sostituendo nella \eqref{eq:somma_discreta} si ottiene
\begin{equation}
    \frac{d\sigma}{d\Omega} = |Z f(\theta,\phi)|^2 \left| \int_{V}n(\mathbf{r}) e^{i\mathbf{q}\cdot\mathbf{r}} d^3r \right|^2
\end{equation}
ed essendo l'ampiezza di diffusione proporzionale alla carica del centro, $Z f(\theta,\phi)$ è l'ampiezza di un unico bersaglio puntiforme di carica $Ze$.

Tuttavia nel nostro caso i bersagli puntiformi sono cariche, quindi possiamo considerare la densità $n(\mathbf{r})$, definita sopra, come la densità di carica normalizzata 
\begin{equation}
    n(\mathbf{r}) = \frac{\rho(\mathbf{r})}{Ze}
\end{equation}
In questo modo 
\begin{equation}
    \int_{V}n(\mathbf{r}) e^{i\mathbf{q}\cdot\mathbf{r}} d^3r = \frac{\int_{V}\rho(\mathbf{r}) e^{i\mathbf{q}\cdot\mathbf{r}} d^3r}{Ze} = \frac{\int_{V}\rho(\mathbf{r}) e^{i\mathbf{q}\cdot\mathbf{r}} d^3r}{\int_V \rho(\mathbf{r}) d^3r}
\end{equation}
Andiamo a chiamare questa quantità fattore di forma\mkeyword{fattore di forma}
\begin{tcolorbox}[colback=black!5!white, colframe=black!45!white, boxrule=0.8pt]
\begin{equation}
    F(\mathbf{q}) = \frac{\int_{V}\rho(\mathbf{r}) e^{i\mathbf{q}\cdot\mathbf{r}} d^3r}{\int_V \rho(\mathbf{r}) d^3r}
\end{equation}
\end{tcolorbox}
In questo modo la sezione d'urto differenziale è esprimibile come 
\begin{equation}
    \frac{d\sigma}{d\Omega} = |Z f(\theta,\phi)|^2 \left| F(\mathbf{q}) \right|^2
\end{equation}
E ricordando che $|Z f(\theta,\phi)|^2$ è la sezione d'urto differenziale di un bersaglio puntiforme di carica $Ze$ si ha
\begin{tcolorbox}[colback=black!5!white, colframe=black!45!white, boxrule=0.8pt]
\begin{equation}
    \frac{d\sigma}{d\Omega} = \frac{d\sigma}{d\Omega}\bigg|_{\text{punt}} \left| F(\mathbf{q}) \right|^2
\end{equation}  
\end{tcolorbox}

\subsection{Il fattore di forma}

Il risultato della sezione precedente separa il problema in due parti indipendenti: $\frac{d\sigma}{d\Omega}\big|_{\text{punt}}$ contiene tutta la dinamica dell'interazione ma nessuna informazione sull'estensione del bersaglio, mentre $F(\mathbf{q})$ contiene tutta la geometria della distribuzione di carica ma nessuna informazione sulla dinamica. Conviene allora studiare separatamente il fattore di forma.

La prima proprietà segue direttamente dalla definizione: per $\mathbf{q} = 0$ tutti i fattori di fase valgono $1$ e si ha
\begin{equation}
    F(0) = \frac{\int_V \rho(\mathbf{r})\,d^3r}{\int_V \rho(\mathbf{r})\,d^3r} = 1
    \label{eq:F_zero}
\end{equation}
Il fattore di forma non misura dunque quanta carica è presente ma soltanto come essa è distribuita.\\

Va sottolineato che $F(\mathbf{q})$ è una quantità misurata e non calcolata. L'osservabile dell'esperimento è la sezione d'urto differenziale, e dividendola per la sezione d'urto puntiforme si ottiene
\begin{equation}
    \left|F(\mathbf{q})\right|^2 = \frac{\left(d\sigma/d\Omega\right)_{\text{mis}}}{\left(d\sigma/d\Omega\right)_{\text{punt}}}
    \label{eq:F_misurato}
\end{equation}
punto per punto in $\mathbf{q}$, senza che la densità $\rho$ sia mai entrata nel procedimento. Gli esempi che seguono percorrono la strada opposta, calcolando $F$ a partire da una $\rho$ assunta: servono da termine di confronto per la curva sperimentale, perché l'antitrasformata diretta non è praticabile in quanto della \eqref{eq:F_misurato} si misura il modulo quadro, quindi si perde la fase di $F$.

\subsubsection*{Caso a simmetria sferica}

Se la distribuzione di carica è a simmetria sferica, $\rho(\mathbf{r}) = \rho(r)$, l'integrale angolare si può svolgere una volta per tutte. Scegliendo l'asse polare lungo $\mathbf{q}$, in modo che $\mathbf{q}\cdot\mathbf{r} = qr\cos\vartheta$,
\begin{equation*}
    \int_V \rho(r)\,e^{i\mathbf{q}\cdot\mathbf{r}}\,d^3r = 2\pi\int_0^{\infty} r^2\rho(r)\int_{-1}^{1} e^{iqr\cos\vartheta}\,d(\cos\vartheta)\,dr = 4\pi\int_0^{\infty} r^2\rho(r)\,\frac{\sin(qr)}{qr}\,dr
\end{equation*}
da cui
\begin{tcolorbox}[colback=black!5!white, colframe=black!45!white, boxrule=0.8pt]
\begin{equation}
    F(q) = \frac{4\pi}{Ze}\int_0^{\infty} r^2\rho(r)\,\frac{\sin(qr)}{qr}\,dr = \frac{4\pi}{Ze}\int_0^{\infty} r^2\rho(r)\,j_0(qr)\,dr
    \label{eq:F_sferico}
\end{equation}
\end{tcolorbox}
dove $j_0(x) = \sin x / x$ è la funzione di Bessel sferica di ordine zero. Si noti che $F$ dipende ora dal solo modulo $q$ e non più dalla direzione di $\mathbf{q}$: la simmetria sferica della sorgente si trasferisce alla sua trasformata.

Per tradurre $q$ in una quantità misurabile ricordiamo che la diffusione è elastica, $|\mathbf{k}_0| = |\mathbf{k}| = k$, per cui
\begin{equation}
    q^2 = \left|\mathbf{k}_0-\mathbf{k}\right|^2 = 2k^2(1-\cos\theta) = 4k^2\sin^2\!\left(\frac{\theta}{2}\right) \implies q = 2k\sin\!\left(\frac{\theta}{2}\right)
    \label{eq:q_theta}
\end{equation}
È questa relazione a collegare la variabile naturale del fattore di forma all'angolo letto dal rivelatore.


\subsubsection*{Guscio sferico}

Consideriamo per primo il caso in cui tutta la carica sia distribuita su una superficie sferica di raggio $R$,
\begin{equation*}
    \rho(r) = \frac{Ze}{4\pi R^2}\,\delta(r-R)
\end{equation*}
La delta rende immediato l'integrale della \eqref{eq:F_sferico} e si ottiene
\begin{equation}
    F_G(q) = \frac{\sin(qR)}{qR}
    \label{eq:F_guscio}
\end{equation}
Il fattore di forma oscilla con inviluppo $1/qR$ e si annulla per $qR = n\pi$. Il primo minimo si ha dunque in
\begin{equation}
    qR = \pi \implies 2k\sin\!\left(\frac{\theta_{\min}}{2}\right) = \frac{\pi}{R}
    \label{eq:minimo_guscio}
\end{equation}
e misurando l'angolo a cui la sezione d'urto si annulla si risale direttamente a $R$.

\subsubsection*{Sfera uniforme}

Consideriamo ora una carica distribuita uniformemente in un volume sferico,
\begin{equation*}
    \rho(r) = \begin{cases} \dfrac{3Ze}{4\pi R^3} & r < R \\[2mm] 0 & r > R \end{cases}
\end{equation*}
L'integrale si svolge per parti,
\begin{equation*}
    F_S(q) = \frac{3}{qR^3}\int_0^{R} r\sin(qr)\,dr = \frac{3}{qR^3}\left[\frac{\sin(qR)}{q^2} - \frac{R\cos(qR)}{q}\right]
\end{equation*}
da cui
\begin{equation}
    F_S(q) = \frac{3}{(qR)^3}\Big[\sin(qR) - qR\cos(qR)\Big]
    \label{eq:F_sfera}
\end{equation}
Gli zeri non sono più equispaziati: imponendo $F_S = 0$ si ottiene la condizione $\tan(qR) = qR$, il cui primo zero non banale è
\begin{equation}
    qR \simeq 4{,}493
    \label{eq:minimo_sfera}
\end{equation}
Il confronto con la \eqref{eq:minimo_guscio} mostra che il coefficiente che converte la posizione del minimo in un raggio dipende dal modello, e ne dipende in modo non che trascurabile.

\subsubsection*{Significato dei minimi}

La presenza di zeri veri in entrambi gli esempi è una conseguenza del bordo netto. Una distribuzione che si interrompe bruscamente minimi completi, se invece il bordo è sfumato i minimi vengono parzialmente riempiti e l'inviluppo cade più rapidamente, mentre una distribuzione ovunque liscia, come una gaussiana, non presenta alcun minimo e decade in modo monotono.

I nuclei reali sono descritti bene dalla distribuzione di Fermi\mkeyword{distribuzione di Fermi},
\begin{equation}
    \rho(r) = \frac{\rho_0}{1+e^{(r-c)/a}}
    \label{eq:woods_saxon}
\end{equation}
nella quale $c$ è il raggio a metà altezza e $a$ regola lo spessore della superficie, i due parametri che avevamo introdotto all'inizio del capitolo come \emph{raggio medio} e \emph{skin thickness}. La posizione dei minimi fissa $c$, il loro riempimento e la caduta dell'inviluppo fissano $a$.

\subsubsection*{Il raggio quadratico medio}

Esiste infine un modo di estrarre un'informazione sulla distribuzione di carica \emph{senza} assumerne la forma. Sviluppando la funzione di Bessel nella \eqref{eq:F_sferico} per $qr \ll 1$,
\begin{equation*}
    \frac{\sin(qr)}{qr} = 1 - \frac{1}{6}q^2r^2 + O(q^4r^4)
\end{equation*}
si ottiene
\begin{equation*}
    F(q) \simeq \frac{4\pi}{Ze}\int_0^{\infty} r^2\rho(r)\,dr - \frac{q^2}{6}\,\frac{4\pi}{Ze}\int_0^{\infty} r^4\rho(r)\,dr = 1 - \frac{1}{6}q^2\langle r^2\rangle
\end{equation*}
avendo riconosciuto nel primo termine la normalizzazione \eqref{eq:F_zero} e avendo definito il \emph{raggio quadratico medio}\mkeyword{raggio quadratico medio}
\begin{equation}
    \langle r^2\rangle = \frac{\int_V r^2\rho(\mathbf{r})\,d^3r}{\int_V \rho(\mathbf{r})\,d^3r}
    \label{eq:r2_def}
\end{equation}
Isolando $\langle r^2\rangle$ si ottiene la relazione cercata
\begin{tcolorbox}[colback=black!5!white, colframe=black!45!white, boxrule=0.8pt]
\begin{equation}
    \langle r^2\rangle = -6\,\frac{d}{dq^2}F(q)\bigg|_{q\to 0}
    \label{eq:r2_derivata}
\end{equation}
\end{tcolorbox}
La pendenza iniziale del fattore di forma misurato, graficato in funzione di $q^2$, fornisce dunque direttamente il raggio quadratico medio della distribuzione di carica. Il risultato non dipende da alcuna ipotesi sulla forma di $\rho$: a piccolo $q$ tutti i modelli collassano sullo stesso andamento, e l'unica informazione che sopravvive è $\langle r^2\rangle$. È per questo motivo che $\sqrt{\langle r^2\rangle}$ costituisce la definizione operativa di raggio nucleare che cercavamo all'inizio del capitolo.